% Generated draft: requires mathematical and source review.
\documentclass[11pt]{article}
\usepackage[margin=25mm]{geometry}
\usepackage{fontspec}
\setmainfont{Latin Modern Roman}
\usepackage{amsmath,amssymb,mathtools,unicode-math}
\setmathfont{Latin Modern Math}
\usepackage{enumitem,needspace,xurl,hyperref}
\setlength{\parindent}{0pt}
\setlength{\parskip}{5pt}
\setlength{\emergencystretch}{3em}
\newcommand{\N}{\mathbb{N}}
\newcommand{\Z}{\mathbb{Z}}
\newcommand{\Q}{\mathbb{Q}}
\newcommand{\R}{\mathbb{R}}
\newcommand{\C}{\mathbb{C}}
\newcommand{\F}{\mathbb{F}}
\newcommand{\A}{\mathbb{A}}
\newcommand{\PP}{\mathbb{P}}
\newcommand{\E}{\mathbb{E}}
\newcommand{\eps}{\varepsilon}
\newcommand{\dd}{\,\mathrm{d}}
\begin{document}
\section*{UnsolvedMath ID: 30004552}
\textbf{Source problem code:} catalogue entry
\subsection{Definitions and mathematical statement}
Let \(R\) denote the theory called Robinson's theory \(R\) in the supplied record. Its axioms and language are not supplied, so its exact presentation remains to be confirmed. Fix a notion of interpretation, and write \(A\preccurlyeq B\) when \(B\) interprets \(A\). Define strict interpretability by
\[
A\lhd B\quad\Longleftrightarrow\quad A\preccurlyeq B\ \text{and}\ B\not\preccurlyeq A.
\]
The record does not specify the permitted interpretations or theory languages; the formulation below is conditional on fixing these conventions.

A theory is recursively axiomatizable if it admits a recursively enumerable set of axioms, and is consistent if it proves no contradiction. A consistent theory \(S\) is incomplete if there is a sentence \(\varphi\) of its language such that \(S\nvdash\varphi\) and \(S\nvdash\neg\varphi\). Define
\[
\begin{aligned}
G1(T)\quad\Longleftrightarrow\quad
&\text{for every recursively axiomatizable consistent theory }S,\\
&T\preccurlyeq S\ \Longrightarrow\ S\text{ is incomplete}.
\end{aligned}
\]
Let
\[
\mathcal D=\{T:T\text{ is recursively axiomatizable and consistent},\ T\lhd R,\ G1(T)\}.
\]
The supplied corrected statement and dated assessment report that \(\mathcal D\neq\varnothing\). This is supplied background, not an independently verified theorem here.

\begin{enumerate}
\item Is \((\mathcal D,\lhd)\) well founded? Explicitly, does every nonempty subset \(\mathcal E\subseteq\mathcal D\) contain a theory \(T\) for which no \(U\in\mathcal E\) satisfies \(U\lhd T\)? Equivalently, in the usual set-theoretic setting, is there no sequence \((T_n)_{n\in\N}\) in \(\mathcal D\) with \(T_{n+1}\lhd T_n\) for every \(n\in\N\)?
\item Does \(G2\) hold for \(R\), in the sense of the supplied corrected statement:
\[
R\nvdash\operatorname{Con}(R)?
\]
Here \(\operatorname{Con}(R)\) is intended to be an arithmetical sentence expressing that there is no \(R\)-proof of a contradiction. The effective presentation of \(R\), proof coding, and internal consistency formula are not specified in the record; this question therefore requires those choices to be fixed before it is fully precise.
\end{enumerate}

The original statement instead asks whether such a theory \(T\lhd R\) exists and defines \(G2\) using the opposite assertion, that a theory proves its own consistency. The supplied verification note says that the report uses nonprovability of consistency, records the existence result below \(R\), and leaves well-foundedness and \(G2\) open. The questions above follow that supplied correction while retaining this formulation conflict explicitly. The dated assessment's claim to answer the ``first question'' concerns the original existence question; it does not supply an answer to well-foundedness.
\subsection{Short English statement}
Consider recursively axiomatizable consistent theories strictly below Robinson's theory \(R\) under interpretation, such that every recursively axiomatizable consistent theory interpreting them is incomplete. Ask whether this class is well founded under strict interpretation, and whether \(R\) fails to prove its own consistency. The supplied record reports existence of such theories, but the interpretation conventions, presentation of \(R\), and consistency sentence require clarification. Its original and corrected formulations disagree about the meaning of \(G2\).
\subsection{Sources}
\begin{itemize}
\item Local UnsolvedMath registry: supplied complete record with permanent ID \(30004552\), exact source problem code \texttt{catalogue entry}, title \emph{Incompleteness and Self-Consistency Below Robinson Arithmetic}, in \emph{Oberwolfach Reports Open Problems}. The proposed year is supplied as 2020; no individual proposer is identified.
\item Supplied report citation: \emph{Mathematical Logic: Proof Theory, Constructive Mathematics} (2021), \emph{Oberwolfach Reports}, pp.\ 1693--1757, DOI \texttt{10.4171/owr/2020/34}. \url{https://doi.org/10.4171/owr/2020/34}. The background describes these reports as documenting workshops and their open-problem discussions.
\item Supplied literature source: Cheng, \emph{Finding the Limit of Incompleteness I}. \url{https://doi.org/10.1017/bsl.2020.9}. The supplied assessment attributes to Cheng the construction of recursively enumerable theories \(U\lhd R\) satisfying \(G1\).
\item Supplied literature source: \emph{Oberwolfach Report 34/2020}. \url{https://doi.org/10.4171/owr/2020/34}.
\item Literature assessment date: \texttt{2026-08-22}. This is supplied dated curation, not an independently verified assessment. It labels the problem partially solved, reports the below-\(R\) existence result, and describes the \(G2\) question for \(R\) as open in the cited literature. The background explicitly requests independent verification of time-sensitive status and cited claims.
\item The registry labels the statement \texttt{corrected\_verified}. Its supplied verification note says the official report was checked, that \(G2\) means nonprovability of consistency, and that well-foundedness and \(G2\) remain open there. This is a report of supplied verification metadata; no source visit or independent verification is asserted here.
\end{itemize}
\end{document}
